\medskip

L'objet de cet exercice est l'étude de l'arrêt d'un chariot sur un manège, à partir du
moment où il entre dans la zone de freinage en fin de parcours.

On note $t$ le temps écoulé, exprimé en seconde, à partir du moment où le chariot arrive sur la zone de freinage.

On modélise la distance parcourue par le chariot dans la zone de freinage, exprimée en mètre, en fonction de $t$, à l'aide d'une fonction notée $d$ définie sur $[0 ~;~+\infty[$.

On a ainsi $d(0) = 0$.

Par ailleurs, on admet que cette fonction $d$ est dérivable sur son ensemble de
définition. On note $d'$ sa fonction dérivée.

\begin{center}\textbf{Partie A} \end{center}

Sur la figure (Fig. 2) ci-dessous, on a tracé dans un repère orthonormé :

\begin{itemize}[label=$\bullet~$]
\item la courbe représentative $\mathcal{C}_d$ de la fonction $d$ ;
\item la tangente $T$ à la courbe $\mathcal{C}_d$ au point A d'abscisse 4,7 ;
\item l'asymptote $\Delta$ à $\mathcal{C}_d$ en $+\infty$.
\end{itemize}

\begin{center}
\psset{unit=0.4cm,arrowsize=2pt 3}
\begin{pspicture}(-1,-2)(17,24)
\psgrid[gridlabels=0pt,subgriddiv=5,gridcolor=gray,subgridcolor=lightgray](0,0)(17,24)
\rput(8.25,-2){$Fig. 2$}
\psplot[plotpoints=2000,linewidth=1.25pt,linecolor=blue]{0}{16.5}{205 9 div 5 3 div x mul 205 9 div add 2.71828 0.6 x mul exp div sub}
\psaxes[linewidth=1.25pt,Dx=10,Dy=10]{->}(0,0)(0,0)(17,24)
\uput[ul](4.7,21){A}\uput[d](16,0){$t$}\uput[ul](7,23){$T$}
\psline(0,16.4)(8,24.3)\psdot(4.7,21)\uput[d](16,22.7){\blue $\mathcal{C}_d$}
\uput[d](1,0){1}\uput[l](0,1){1}
\psline[linewidth=1.25pt,linestyle=dashed](0,22.778)(17,22.778)\uput[u](2,22.778){$\Delta$}
\end{pspicture}
\end{center}


\emph{Dans cette partie, aucune justification n'est attendue.}

\medskip

Avec la précision que permet le graphique, répondre aux questions ci-dessous.

D'après ce modèle :

\medskip

\begin{enumerate}
\item Au bout de combien de temps le chariot aura-t-il parcouru $15$~m dans la zone de freinage ?
\item Quelle longueur minimale doit-être prévue pour la zone de freinage ?
\item Que vaut $d'(4,7)$ ? Interpréter ce résultat dans le contexte de l'exercice.
\end{enumerate}

\begin{center}\textbf{Partie B} \end{center}

On rappelle que $t$ désigne le temps écoulé, en seconde, à partir du moment où le chariot arrive sur la zone de freinage.

On modélise la vitesse instantanée du chariot, en mètre par seconde $\left(\text{m}.\text{s}^{-1}\right)$, en fonction de $t$, par une fonction $v$ définie sur $[0~;~+\infty[$.

On admet que :

\begin{itemize}[label=$\bullet~$]
\item la fonction $v$ est dérivable sur son ensemble de définition, et on note $v'$ sa
fonction dérivée ;
\item la fonction $v$ est une solution de l'équation différentielle 
\[(E) :\quad y' + 0,6y = \e^{-0,6t},\]

où $y$ est une fonction inconnue et où $y'$ est la fonction dérivée de $y$.

On précise de plus que, lors de son arrivée sur la zone de freinage, la vitesse du chariot est égale à $12$ m.s$^{-1}$, c'est-à-dire $v(0) = 12$.

\end{itemize}

\medskip

\begin{enumerate}
\item 
	\begin{enumerate}
		\item On considère l'équation différentielle 
		\[(E') :\quad y' + 0,6y = 0.\]
		
Déterminer les solutions de l'équation différentielle $(E')$ sur $[0~;~+\infty[$.
		\item Soit $g$ la fonction définie sur $[0~;~+\infty[$ par $g(t) = t\e^{-0,6t}$.
		
Vérifier que la fonction $g$ est une solution de l'équation différentielle $(E)$.
		\item En déduire les solutions de l'équation différentielle $(E)$ sur $[0~;~+\infty[$.
		\item En déduire que pour tout réel $t$ appartenant à l'intervalle $[0~;~+\infty[$, on a :
\[v(t) = (12+t)\e^{-0,6t}.\]
	\end{enumerate}
\item Dans cette question, on étudie la fonction $v$ sur $[0~;~+\infty[$.
	\begin{enumerate}
		\item Montrer que pour tout réel $t \in [0~;~ +\infty[,\: v'(t) = (-6,2 - 0,6t)\e^{-0,6t}$.
		\item En admettant que:
		
\[v(t) = 12\e^{-0,6t} + \dfrac{1}{0,6} \times  \dfrac{0,6t}{\e^{0,6t}},\]
déterminer la limite de $v$ en $+\infty$.
		\item Étudier le sens de variation de la fonction $v$ et dresser son tableau de
variation complet. Justifier.
		\item Montrer que l'équation $v(t) = 1$ admet une solution unique $\alpha$, dont on
donnera une valeur approchée au dixième.
	\end{enumerate}
\item Lorsque la vitesse du chariot est inférieure ou égale à 1 mètre par seconde, un système mécanique se déclenche permettant son arrêt complet.

Déterminer au bout de combien de temps ce système entre en action. Justifier.
\end{enumerate}

\begin{center}\textbf{Partie C} \end{center}

On rappelle que pour tout réel $t$ appartenant à l'intervalle $[0~;~ +\infty[$ :

\[v(t) = (12 + t)\e^{-0,6t}.\]

On admet que pour tout réel $t$ dans l'intervalle $[0~;~ +\infty[$ :

\[d(t) = \displaystyle\int_0^t v(x)\:\text{d}x.\]

\smallskip

\begin{enumerate}
\item À l'aide d'une intégration par parties, montrer que la distance parcourue par le
chariot entre les instants 0 et $t$ est donnée par : 

\[d(t)= \e^{-0,6t} \left(- \dfrac53 t - \dfrac{205}{9}\right) + \dfrac{205}{9}.\]

\item On rappelle que le dispositif d'arrêt se déclenche lorsque la vitesse du chariot est inférieure ou égale à 1 mètre par seconde.

Déterminer, selon ce modèle, une valeur approchée au centième de la distance parcourue par le chariot dans la zone de freinage avant le déclenchement de ce dispositif.
\end{enumerate}
